\section{Levantamientos locales, filas y control uniforme}
\label{sec:levantamientos-uniformes}

\fuente{Continuación de la Nota~\ref{sec:theta-integral}.
A93, §9, pp.~62--64, para condiciones locales, finitud y normalización.
El cotejo de los productores locales conserva el espejo
\texttt{CocoSgt/riemann-hypothesis-lean} en la revisión
\texttt{7cbf7a0c32fca1badf05931879a37c3f8a2b03b4}.
Se distinguen las pruebas desarrolladas aquí, los resultados previos
que se aplican y las entradas de teoría de periodos recibidas.}

La Nota~\ref{sec:theta-integral} construyó los cocientes integrales
que miden el defecto de separación y dejó escrito el presupuesto
\eqref{eq:ti-local-display}. Esta nota desarrolla los mapas locales
que permiten alimentar ese presupuesto. El primer objetivo es
pasar de una clase levantada a un cociclo con la igualdad exacta
que exigen las relaciones de columnas. El segundo es identificar
el cociente de una fila y estimar su longitud sin perder su torsión.
El tercero es cuantificar los núcleos de invariantes y de
restricción cuando cambian los datos aritméticos.

La distinción que guía las pruebas es sencilla: disponer de un
objeto finito para cada miembro de una familia no ofrece todavía
una constante común. Obtendremos certificados concretos de esa
constante: exponentes de levantamiento, valoraciones de menores,
cortes de caracteres y grados de extensiones. La aplicación a la
familia theta conserva las hipótesis que realizan esos certificados.

En toda la nota $K/\Q_\ell$ es finita, $\mathcal O$ es su anillo
de enteros, $\varpi$ un uniformizante y
$e_{\mathcal O}=v_\varpi(\ell)$. La longitud puede tomar el valor
infinito; solo escribimos
$\lambda_{\mathcal O}(M)=\length_{\mathcal O}(M)/e_{\mathcal O}$
como número finito después de justificar su finitud. Los cuerpos
aritméticos se denotan por $E$ o $L$, para distinguirlos del
cuerpo de coeficientes $K$.

\subsection{Qué significa levantar una clase local con coste acotado}
\label{subsec:pl-levantar}

Sea $\mathcal O$ un anillo de valoración discreta de característica
mixta $(0,\ell)$, con cuerpo de fracciones $K$, uniformizante
$\varpi$ y $e_{\mathcal O}=v_\varpi(\ell)$. Para un grupo $D$ que
actúa $\mathcal O$-linealmente en $M$, escribimos
\[
 \begin{gathered}
 Z^1(D,M)=\{z:D\to M:z(gh)=z(g)+g z(h)\},\\
 d_Mm(g)=gm-m,\qquad H^1(D,M)=Z^1(D,M)/d_MM.
 \end{gathered}
\]
En esta primera construcción la cohomología es algebraica. La versión
continua se indicará expresamente cuando intervenga la topología.

Sean $p:M_1\to M_2$ $\mathcal O$-lineal y equivariante, $g=H^1(p)$ y
$H_{f,i}\subseteq H^1(D,M_i)$ dos submódulos que fijan condiciones
locales. Para un
entero $e\ge0$, la afirmación
\begin{equation}
 \ell^e H_{f,2}\subseteq g(H_{f,1})
 \label{eq:pl-lifto}
\end{equation}
acota el exponente del conúcleo local. No afirma que un representante
arbitrario de la clase levantada se envíe exactamente en
$\ell^e\kappa$. Para obtener esa igualdad hay que corregir un borde.

\begin{proposicion}[Levantamiento exacto de cociclos]
\label{prop:pl-cociclo-exacto}
Supongamos que $p$ es sobreyectiva y que se cumple
\eqref{eq:pl-lifto}. Para cada $\kappa\in Z^1(D,M_2)$ con
$[\kappa]\in H_{f,2}$ existe $K_\kappa\in Z^1(D,M_1)$ tal que
\[
 [K_\kappa]\in H_{f,1},\qquad
 p(K_\kappa(g))=\ell^e\kappa(g)\quad(g\in D).
\]
La corrección no aumenta el exponente $e$.
\end{proposicion}
\begin{proof}
Por \eqref{eq:pl-lifto} elegimos una clase $[K_0]\in H_{f,1}$ con
$[pK_0]=\ell^e[\kappa]$. Existe, por tanto, $v\in M_2$ tal que
$pK_0-\ell^e\kappa=d_{M_2}v$. Elegimos $w\in M_1$ con $p(w)=v$
y ponemos $K_\kappa=K_0-d_{M_1}w$. La equivariancia da
$p(d_{M_1}w)=d_{M_2}v$, de modo que $pK_\kappa=\ell^e\kappa$.
Restar un borde preserva la clase de cohomología y, por tanto, su
pertenencia a $H_{f,1}$.
\end{proof}

Esta es la prueba de \texttt{LocLift.liftCoc}, líneas 13--36, y no
solo una traducción de su enunciado. La sobreyectividad sobre
\emph{los módulos integrales} permite elegir $w$ sin denominador.
En la versión continua también se necesitan un representante $K_0$
continuo y continuidad del mapa orbital $g\mapsto gw$; bajo esas
hipótesis la misma corrección es continua.

\subsection{Las dos evaluaciones de las columnas}
\label{subsec:pl-columnas}

Sea $A$ una $\mathcal O$-álgebra, y sea la acción de $D$ en los
módulos considerados la multiplicación por $u(g)$. En la aplicación
$u(g)=\chi(g)^{-1}r(g)$, con el carácter tomando valores centrales.
Sea $P$ un $A$-módulo con un mapa $A$-lineal
$\iota:P\to A$. Supongamos dada una evaluación sobreyectiva
\[
 \operatorname{ev}_{v_1}:P\longrightarrow L,
 \qquad p\longmapsto\iota(p)v_1,
\]
que es equivariante para las acciones usadas. Para una familia finita
$\kappa_a\in Z^1(D,L)$ cuyas clases satisfacen la condición local,
\eqref{eq:pl-lifto} y la proposición anterior producen cociclos
$z_a:D\to P$. Entonces
\begin{equation}
 K_a(g):=\iota(z_a(g)),\qquad
 K_a(g)v_1=\ell^e\kappa_a(g),\qquad [z_a]\in H_{f,P}.
 \label{eq:pl-colk}
\end{equation}
En particular, si la condición finita es functorial para los mapas
$P\to V$, $p\mapsto\iota(p)s$, también la función
$g\mapsto K_a(g)s$ satisface esa condición, para cada $s\in V$.
Esto es lo que utilizan \texttt{colK} y \texttt{colFun}, líneas
38--49 y 68--77 de \texttt{LocLift.lean}.

La segunda evaluación tiene como objetivo el ideal
$\delta_{\mathcal O}\mathcal O$, con acción trivial. Supongamos que
\[
 \operatorname{ev}_E:P\twoheadrightarrow\delta_{\mathcal O}\mathcal O,
 \quad \iota(P)\subseteq AE,
 \quad \iota(p)m=\operatorname{ev}_E(p)m\quad(p\in P,m\in M).
\]
Sea $g\mapsto\delta_{\mathcal O}u_0(g)$ un cociclo en ese ideal y
supongamos que su clase satisface la condición local objetivo.
Supongamos además que $\operatorname{ev}_E$ es equivariante para
la acción trivial del objetivo y satisface \eqref{eq:pl-lifto}
con sus condiciones locales correspondientes y el mismo exponente $e$.
Aplicando la misma proposición obtenemos $z_0:D\to P$ y
\begin{equation}
 k_0(g)=\iota(z_0(g))\in AE,\qquad
 k_0(g)m=\ell^e\delta_{\mathcal O}u_0(g)m.
 \label{eq:pl-col-e}
\end{equation}
Aquí $\delta_{\mathcal O}$ forma parte del cociclo de entrada. La
corrección de bordes no introduce una segunda potencia de
$\delta_{\mathcal O}$.

Cuando $D$ es un subgrupo de un grupo global $G$, el código define
$K_a$ y $k_0$ en $G\setminus D$ por cero. Esa elección prolonga
\emph{funciones}; las identidades de cociclo y las condiciones
locales se utilizan en $D$. Sin hipótesis adicionales sobre $D$ no
se deduce que la prolongación sea un cociclo global ni que sea
continua. Los cociclos que sí se han construido son $z_a$ y $z_0$
en el grupo local.

\subsection{Por qué basta una suma finita de representaciones matriciales}
\label{subsec:pl-seleccion}

Para aplicar un teorema de levantamiento a $P$ hay que situar
$K\otimes_{\mathcal O}P$ en una clase de representaciones a la que
ese teorema se aplique. La evaluación en la familia de matrices
ofrece una construcción concreta.

\begin{lema}[Selección finita de columnas]
\label{lem:pl-seleccion}
Sea $V$ un espacio vectorial de dimensión finita sobre $K$, con
acción de $D$. Supongamos dado un mapa lineal inyectivo
\[
 \Lambda:V\longrightarrow\prod_{x\in X}\operatorname{Mat}_d(K)
\]
tal que $\Lambda(gv)_x=U_x(g)\Lambda(v)_x$ para una familia de
representaciones matriciales $U_x$. Existen
$m\leq\dim_KV$ y pares $(x_a,b_a)$, $1\leq a\leq m$, tales que
\[
 j:V\hookrightarrow\bigoplus_{a=1}^m K^d,
 \qquad j(v)_a=\Lambda(v)_{x_a}e_{b_a},
\]
es inyectiva y equivariante, con acción $U_{x_a}$ en el sumando $a$.
\end{lema}
\begin{proof}
Las formas lineales que extraen todas las entradas de todas las
matrices $\Lambda(v)_x$ tienen aniquilador común cero, pues
$\Lambda$ es inyectiva. Por dualidad en dimensión finita, generan
$V^*$. Elegimos entre ellas una base de $V^*$, con $m=\dim_KV$
elementos (y la familia vacía si $V=0$). A cada forma seleccionada
le corresponde una columna $(x_a,b_a)$. Si todas esas columnas
se anulan en $v$, se anulan las formas de la base, luego $v=0$.
La identidad de equivariancia resulta de que extraer una columna
conmuta con la multiplicación matricial a la izquierda.
\end{proof}

La selección finita corresponde a \texttt{fin\_sel}, líneas
593--619 de \texttt{FiniteClassesLocalEq.lean}; la extensión de la
evaluación integral a $K\otimes P$ y la extracción de columnas
aparecen en \texttt{embAt\_Dd}, líneas 708--776. La cota
$m\leq\dim_KV$ se ha explicitado aquí mediante la prueba anterior;
no se necesita una cota uniforme de $\dim_KV$ para la mera
existencia de la inmersión. Para usar un mismo exponente en toda la
familia se necesita, en cambio, que la hipótesis de levantamiento
sea uniforme respecto del número de sumandos. Ese cuantificador
está incluido en \texttt{LiftAt}: se permiten cualesquiera
subcocientes de sumas \emph{finitas} de la familia prescrita, tanto
en la fuente como en el objetivo.

Los productores \texttt{liftCol\_eve\_Dd} y
\texttt{liftCol\_evE\_Dd}, líneas 798--841 del mismo archivo,
usan esa inmersión, la continuidad de la familia y la entrada
\texttt{LiftExponentAtPlaces} para obtener \eqref{eq:pl-lifto}
con el mismo exponente local. La operación que transporta
\texttt{LiftF} a \texttt{LiftO}, líneas 178--237, conserva el
exponente: la discrepancia es un borde integral que desaparece al
pasar a cohomología. Se mantienen las hipótesis de retículos
finitos y libres, evaluación sobreyectiva, compatibilidad de
coeficientes y entradas locales requeridas por el productor.

\subsection{El cociente de una fila, antes de estimar su longitud}
\label{subsec:pl-cociente-fila}

Ahora sí usamos cohomología continua. Sea $D$ topológico y sean
$T_i\subset V_i$ retículos estables en representaciones de
dimensión finita sobre $K$. Suponemos continuos los mapas
orbitales y las inclusiones de los retículos, con las topologías
compatibles. Sea $f:T_1\to T_2$ equivarante, extendido a un mapa
continuo $f_K:V_1\to V_2$.

Sea $Z_f^1(D,V_i)$ una condición finita sobre cociclos continuos,
que sea un $K$-subespacio, contenga $dV_i$ y sea preservada por
$f_K$. En el modelo de periodos leído, su definición es
\begin{equation}
 Z_f^1(D,V)=\{z\in Z^1_{\rm cont}(D,V):
   \exists b\in V\otimes_{\mathbb Q_\ell}B,\quad
   z(g)\otimes1=gb-b\text{ para todo }g\in D\}.
 \label{eq:pl-zf}
\end{equation}
La acción en $V\otimes_{\mathbb Q_\ell}B$ es la acción diagonal prescrita por
los datos de periodos. Esta definición corresponde a
\texttt{finiteCocycles}, líneas 190--202 de
\texttt{WeaklyAdmissibleFiltrations/}\allowbreak\texttt{Statement.lean}.
El adjetivo ``finita'' describe una condición aritmética: por sí
solo no afirma finitud cardinal ni longitud finita.

Ponemos
\[
 Z_f^1(D,T_i)=\{z\in Z^1_{\rm cont}(D,T_i):z_K\in Z_f^1(D,V_i)\},
 \qquad H_f^1(D,T_i)=\operatorname{im}(Z_f^1(D,T_i)\to H^1(D,T_i)).
\]
Definimos también $H_f^1(D,V_2)$ como el $K$-espacio generado por
las clases finitas racionales, y
\begin{equation}
 F_s=\{\xi\in H^1(D,T_2):\xi_K\in H_f^1(D,V_2)\},\qquad
 I=f_*H_f^1(D,T_1).
 \label{eq:pl-fs}
\end{equation}
Las propiedades anteriores implican $I\subseteq F_s$.

\begin{proposicion}[Presentación exacta del defecto de fila]
\label{prop:pl-presentacion-fila}
Con estas hipótesis, para todo cociclo integral continuo $z$,
\[
 [z]\in F_s\quad\Longleftrightarrow\quad z_K\in Z_f^1(D,V_2).
\]
Además, el mapa que toma la clase de un cociclo induce un
isomorfismo de $\mathcal O$-módulos
\begin{equation}
 \frac{F_s}{I}\ \simeq
 \frac{Z_f^1(D,T_2)}{fZ_f^1(D,T_1)+dT_2}.
 \label{eq:pl-iso-fila}
\end{equation}
Por tanto, las dos presentaciones tienen exactamente la misma
longitud, finita o infinita.
\end{proposicion}
\begin{proof}
La implicación de derecha a izquierda es inmediata. Si la clase
racional de $z$ pertenece al espacio generado por las clases
finitas, tomamos una combinación finita de representantes finitos.
Equivalentemente, tras despejar sus denominadores, existen
$0\ne t\in\mathcal O$, $c\in Z_f^1(D,V_2)$ y $v\in V_2$ con
$tz_K=c+d v$. Como $Z_f^1(D,V_2)$ es un $K$-subespacio y contiene
$dV_2$, se sigue que $z_K=t^{-1}(c+d v)$ también pertenece a él.
La elección del signo del borde no altera el argumento.

En consecuencia, $Z_f^1(D,T_2)\to F_s$ es sobreyectiva y tiene
núcleo $dT_2$. La preimagen de $I$ consiste exactamente en los
cociclos $f(c)+dt$, con $c\in Z_f^1(D,T_1)$ y $t\in T_2$.
La functorialidad y la inclusión de los bordes aseguran que esa
suma ya está contenida en $Z_f^1(D,T_2)$. El teorema de isomorfía
da \eqref{eq:pl-iso-fila}.
\end{proof}

Para condiciones locales arbitrarias $F_i\subseteq H^1(D,T_i)$,
si $Z_{F_i}$ denota su preimagen en los cociclos, la fórmula que
no presupone functorialidad es
\[
 \frac{F_2}{F_2\cap f_*F_1}\simeq
 \frac{Z_{F_2}}{Z_{F_2}\cap(fZ_{F_1}+dT_2)}.
\]
La simplificación del denominador en \eqref{eq:pl-iso-fila}
utiliza precisamente la inclusión $f_*F_1\subseteq F_2$.

El denominador es $dT_2$, no el conjunto más grande de los bordes
racionales que casualmente tengan valores integrales. Reemplazarlo
por ese conjunto puede eliminar parte del defecto. En particular,
$F_s$ puede tener torsión procedente del núcleo de
$H^1(D,T_2)\to H^1(D,V_2)$.

En el código, $\operatorname{col}(F)$ se define como la longitud
del cociente de la derecha (formalmente se intersecta el
denominador con el numerador). Las funciones
\texttt{FsS}, \texttt{ImS} y \texttt{rowLenS} emplean el cociente
de la izquierda. La prueba de \texttt{LiftMapEquivariant.mem\_FsG},
líneas 91--134, efectúa el paso por el espacio racional;
\texttt{im\_iff}, líneas 208--280, conserva los bordes integrales;
y \texttt{len\_eq\_col}, líneas 282--302, identifica las longitudes.

\subsection{La topología que permite comparar los dos modelos}
\label{subsec:pl-topologia}

Los grupos locales y los coeficientes de la definición de periodos
pueden estar presentados de modo distinto a los grupos y retículos
de la fila. Algebraicamente, una isomorfía de grupos y una
isomorfía de anillos transportan cociclos, bordes y sus cocientes.
La continuidad del cociclo transportado requiere una comprobación
adicional. En \texttt{rowLenD\_eq} y \texttt{len\_eq\_col} se
usan dos condiciones de reflexión de continuidad:

\begin{enumerate}
\item una función con valores racionales es continua si lo es su
      transporte al modelo local de periodos;
\item una función con valores en el retículo fuente es continua
      si lo es su racionalización transportada.
\end{enumerate}

Estas condiciones se satisfacen, por ejemplo, cuando los cambios
de grupo y espacio racional son homeomorfismos y el retículo
lleva la topología inducida con la inclusión correspondiente.
El código las retiene como las entradas \texttt{RowTop}; no deben
suprimirse de un argumento que parta de cociclos algebraicos.
La inclusión de los bordes racionales en \eqref{eq:pl-zf} utiliza
también la continuidad orbital: para $dv$ se toma $v\otimes1$
como testigo del periodo.

Una isomorfía semilineal para una \emph{isomorfía de anillos}
preserva la longitud: transporta biyectivamente las cadenas de
submódulos, con sus inclusiones estrictas y sus factores simples.
Por ello el transporte de \texttt{len\_eq\_col} no tiene coste
de ramificación. El factor
$\operatorname{length}_{\mathcal O}(\mathcal O/\ell\mathcal O)
=e_{\mathcal O}$ que aparecerá en la cota de fila procede de su
normalización, no de este cambio de presentación.

\subsection{Cuánto cuesta el conúcleo: rango y torsión}
\label{subsec:pl-coste}

\begin{proposicion}[Cota de longitud a partir de LiftO]
\label{prop:pl-coste}
Sea $F_s$ un $\mathcal O$-módulo finitamente generado, $I\subseteq
F_s$ y supongamos $\ell^e F_s\subseteq I$. Escribamos
\[
 F_s\simeq\mathcal O^r\oplus
           \bigoplus_{j=1}^t\mathcal O/\varpi^{a_j}\mathcal O,
 \qquad a_j\geq1.
\]
Entonces
\begin{equation}
\begin{aligned}
 \operatorname{length}_{\mathcal O}(F_s/I)
 &\leq r e e_{\mathcal O}+\sum_{j=1}^t\min(a_j,e e_{\mathcal O})
 \\
 &\leq r e e_{\mathcal O}
      +\operatorname{length}_{\mathcal O}(F_{s,\mathrm{tors}}).
\end{aligned}
 \label{eq:pl-coste-fino}
\end{equation}
\end{proposicion}
\begin{proof}
La inclusión $\ell^eF_s\subseteq I$ da una sobreyección
$F_s/\ell^eF_s\twoheadrightarrow F_s/I$. Como
$\ell^e$ es una unidad por $\varpi^{e e_{\mathcal O}}$, cada
sumando libre contribuye $e e_{\mathcal O}$ a la longitud del
primer cociente; cada sumando de torsión contribuye
$\min(a_j,e e_{\mathcal O})$. La aditividad de la longitud da
la primera desigualdad de \eqref{eq:pl-coste-fino}; la segunda
usa $\min(a_j,e e_{\mathcal O})\le a_j$.
\end{proof}

La segunda desigualdad es el contenido de
\texttt{TorsionLengthCounts/RowCol.rowColength}, líneas 83--143,
aplicado por \texttt{LocLift.rowLift}, líneas 79--87. La primera
es una precisión mediante factores invariantes que hemos probado
aquí. Si la imagen inicial $g(H_{f,1})$ no se ha identificado aún
como submódulo de $F_s$, se toma
$I=F_s\cap g(H_{f,1})$; la hipótesis \texttt{LiftO} sigue dando
$\ell^eF_s\subseteq I$.

El término de torsión tiene contenido. Para
$F_s=\mathcal O/\varpi^a\mathcal O$, $I=0$ y
$e e_{\mathcal O}\geq a$, se cumple la hipótesis aunque
$\operatorname{rank}_{\mathcal O}F_s=0$ y
$\operatorname{length}_{\mathcal O}(F_s/I)=a$. Una cota que use
solo el rango daría erróneamente cero. La eliminación de este
término requiere una hipótesis adicional de ausencia de torsión
o una estimación específica de ella.

\subsection{De la estimación local de periodos a TieRow}
\label{subsec:pl-tie}

La presentación \eqref{eq:pl-iso-fila} permite usar una cota
obtenida directamente en el modelo de periodos. En la revisión
cotejada, \texttt{RowAt} exige lo siguiente para cada mapa
sobreyectivo de retículos $F:T_1\to T_2$ de la clase considerada:
el retículo objetivo debe identificarse con el retículo matricial
prescrito; la representación fuente debe ser una
\emph{subrepresentación} de una suma finita de la familia dual;
y los mapas orbitales deben ser continuos. Su conclusión
\texttt{ColLe} consiste en la existencia de $m,k\in\mathbb N$ con
\begin{equation}
 \operatorname{col}(F)=m,\qquad
 \operatorname{length}_{\mathcal O}(\mathcal O/\ell\mathcal O)=k,
 \qquad m\leq b\,k,
 \label{eq:pl-colle}
\end{equation}
donde $b\in\mathbb Q$ es la cota normalizada de la fila. Incluye
la finitud de las dos longitudes; no es una desigualdad obtenida
asignando un número a una longitud infinita.

Las hipótesis de retículo, acción, topología y comparación de
condiciones locales identifican
$\operatorname{rowLen}(n)=\operatorname{col}(F_n)$. La condición
\texttt{SelfGood} permite usar el modelo de coeficientes del
propio punto, para el que $k=e_{\mathcal O}$. Aplicando
\eqref{eq:pl-colle} con $b=\operatorname{lenRow}(w_n)$ se obtiene
\begin{equation}
 \operatorname{rowLen}(n)<\infty,\qquad
 \operatorname{rowLen}(n)\leq
 e_{\mathcal O}\operatorname{lenRow}(w_n)
 \quad(n\in\operatorname{rowsN}).
 \label{eq:pl-tie}
\end{equation}
La ruta cotejada es
\texttt{rowTie\_disp} (líneas 102--157 de
\texttt{FiniteClassesRowEqGeneric.lean}) y
\texttt{BKaClause.tie\_clause} (líneas 60--77).
La igualdad de modelos es el paso que hace aplicable la cota;
la hipótesis \texttt{LengthBoundAtPlaces} proporciona la
estimación local que alimenta esa ruta.

No se han identificado \texttt{LiftO} y \texttt{TieRow} como una
misma entrada. Un levantamiento de exponente acotado proporciona
la cota \eqref{eq:pl-coste-fino} si se controlan rango y torsión;
la entrada \texttt{RowAt} proporciona directamente una cota
normalizada para el cociente particular, con sus hipótesis de
familia dual y retículo matricial.

\subsection{Proveedor por ventanas y frontera de esta reconstrucción}
\label{subsec:pl-ventanas}

Hay un productor concreto del exponente local, aunque requiere
entradas de teoría de periodos. El teorema
\texttt{lifting\_condition\_of\_windows}, líneas 70--88 de
\texttt{LiftingConditionOfWindows.lean}, trabaja para $\ell\ne2$.
Bajo su teorema de levantamiento para todos los casos, las
compatibilidades de Fontaine, la admisibilidad indicada, la
continuidad de la familia, la condición no ramificada de las bases
de columna y una ventana común de saltos $[-M,h]$, basta elegir
\begin{equation}
 e'_w\ \geq\ C_A\left(2+\frac{M+h}{\ell}
                       +\log_\ell(M+h+1)\right),
 \qquad C_A\geq0,
 \label{eq:pl-ventana}
\end{equation}
para obtener \texttt{LiftExponentAtPlaces}. El cuantificador de
la ventana abarca los lugares de columna, las extensiones de
coeficientes declaradas buenas, los miembros de la familia y los
encajes que indexan los saltos.

La reconstrucción anterior prueba los pasos posteriores: la
selección finita de representaciones de columnas, el transporte a
\texttt{LiftO}, la corrección exacta de bordes, la presentación del
cociente de filas y su coste de longitud. El teorema general de
levantamiento, las entradas cristalinas y la identificación del
retículo matricial objetivo siguen siendo hipótesis de los
productores leídos. El archivo de ventanas no prueba por sí solo
esas entradas ni el caso $\ell=2$. Por su parte,
\texttt{RowTieInOf.rowTieIn\_of}, líneas 30--55, ensambla las
propiedades estructurales de la fila, pero recibe expresamente
\texttt{RowMat} y la continuidad orbital necesaria para los bordes.
La presente lectura no sustituye esos datos por conclusiones.

\begin{proposicion}[Una elección común para los exponentes de columna]
\label{prop:lu-ventana-comun}
Supongamos que las hipótesis del productor anterior se realizan
para cada miembro de una familia, con $C_*,W_*\ge0$,
$\ell\ge3$, $C_A\le C_*$ y
$0\le M+h\le W_*$. Entonces puede elegirse en todas sus columnas
el mismo exponente entero
\begin{equation}
 E_*:=\left\lceil C_*\left(2+\frac{W_*}{3}
                       +\log_3(W_*+1)\right)\right\rceil.
 \label{eq:lu-exponente-comun}
\end{equation}
Esta elección no exige acotar el número de sumandos finitos
admitidos por \texttt{LiftAt}.
\end{proposicion}
\begin{proof}
Para $W\ge0$ y $\ell\ge3$ se tiene
$W/\ell\le W/3$ y
$\log_\ell(W+1)\le\log_3(W+1)$. Ambas funciones del lado derecho
son crecientes en $W$. Como $C_A\ge0$, sustituir $W$ por $W_*$
y $C_A$ por $C_*$ hace que \eqref{eq:pl-ventana} se cumpla con
$e'_w=E_*$. El productor cuantifica sobre todas las sumas finitas
de la clase prescrita, de modo que no introduce un factor por
su número de sumandos.
\end{proof}

La proposición fija una elección posible de $e'_w$; no deduce
que cualquier elección mayor ya esté acotada por $E_*$. Tampoco
construye la ventana común, ni incorpora el primo $2$, ni da un
exponente para la familia dual de las filas. Estas son entradas
distintas que deben verificarse en su lugar.

\subsection{Núcleos de columnas e invariantes de la acción torcida}
\label{subsec:lu-h0-fuente}

Para una representación con acción $r$ y un carácter $\chi$, el
operador de columna tiene la fórmula
\begin{equation}
 D_M(m)(g)=\chi(g)^{-1}r(g)m-m.
 \label{eq:lu-d-columna}
\end{equation}
Por definición, $\ker D_M=M^{D,u}$, donde la acción torcida es
$u(g)=\chi(g)^{-1}r(g)$. En el código cotejado, \texttt{colD}
actúa en un espacio racional y su núcleo se mide por dimensión;
\texttt{colDO} actúa en el módulo divisible correspondiente y su
núcleo se mide por longitud. Una dimensión racional finita no
proporciona por sí sola una longitud torsionaria finita.

El archivo \texttt{KernelLengthHigherColumns.lean} prueba dos
comparaciones. Para el bloque $0$, la aplicación \texttt{toKOG}
es inyectiva y equivarante bajo \texttt{RingDef} y \texttt{NormAt}
(líneas 93--135). Restringirla a los invariantes da
\begin{equation}
 h_{0,w}\le\length_{\mathcal O}H^0(E_w,W_{K/\mathcal O}).
 \label{eq:lu-h0-bloque}
\end{equation}
La prueba usa las entradas de las matrices, la identificación de
coeficientes y la acción; una igualdad de rangos no la reemplaza.
Para $c\ge2$, la equivalencia importada
\texttt{blkTw\_of\_def} identifica el bloque con
$K/\mathcal O(-(c-1))$. La misma monotonía da una cota por
\begin{equation}
 H_0:=\max_{\substack{w\in\mathcal C\\1\le e<t}}
       \length_{\mathcal O}H^0(E_w,K/\mathcal O(-e)),
 \label{eq:lu-h0-def}
\end{equation}
con máximo cero si el conjunto de índices es vacío. El bloque
$c=1$ tiene torcido cero y queda fuera de esta fórmula.

Estas comparaciones se demuestran en $\N\cup\{\infty\}$.
La estructura \texttt{H0.FinCore}, líneas 66--72 del mismo archivo,
recibe la finitud de los invariantes del bloque automorfo en su
clase admisible y la finitud de los caracteres negativos en los
lugares sobre $\ell$. No contiene una constante numérica uniforme.
El ensamblaje \texttt{theorem\_of\_display\_and\_g2}, líneas
185--195, utiliza esa entrada. Después se justifica la finitud
del supremo antes de convertir longitudes a naturales y dividir
por $e_{\mathcal O}$. La lectura localizada no afirma la ausencia
de otros productores fuera de los archivos revisados.

\subsection{Un sistema finito detecta todos los invariantes}
\label{subsec:lu-detectores}

Sea $G$ un grupo profinito y sea
$\rho:G\to\operatorname{GL}_d(\mathcal O)$ continua. Escribimos
$A_d=(K/\mathcal O)^d$ con la acción inducida y topología discreta.
Esta acción es continua: todo elemento pertenece a
$A_d[\varpi^n]$ para algún $n$, y su estabilizador contiene el
núcleo abierto de la reducción de $\rho$ módulo $\varpi^n$.
Toda cohomología de $A_d$ en lo que sigue usa cocadenas continuas.

\begin{lema}[A lo sumo $d$ matrices para la invariancia]
\label{lem:lu-detectores}
Sea $R\subseteq\mathcal O^{1\times d}$ el submódulo generado por
todas las filas de todas las matrices $\rho(g)-I_d$, con $g\in G$,
y sea $r=\operatorname{rank}_{\mathcal O}R$. Existen $r$ filas
originales que forman una base de $R$. En particular, pueden
elegirse $s\le r\le d$ elementos $g_1,\ldots,g_s$ tales que,
para la matriz apilada
\begin{equation}
 B=\begin{pmatrix}\rho(g_1)-I_d\\ \vdots\\\rho(g_s)-I_d\end{pmatrix},
 \qquad A_d^G=\ker(B:A_d\to A_{sd}).
 \label{eq:lu-detector}
\end{equation}
Para $r=0$ la familia es vacía y $A_d^G=A_d$.
\end{lema}
\begin{proof}
El anillo $\mathcal O$ es noetheriano y $R$ es un submódulo de
un módulo libre finito. Por tanto $R$ es libre de rango $r$.
Las imágenes de las filas originales generan el espacio
$R/\varpi R$, de dimensión $r$ sobre $k=\mathcal O/\varpi$.
Elegimos $r$ de ellas cuyas imágenes sean una base. Si $R_0$
es el submódulo que generan, entonces
$R=R_0+\varpi R$. El módulo finitamente generado $R/R_0$
satisface $\varpi(R/R_0)=R/R_0$; Nakayama da $R/R_0=0$.
También se ve directamente con la clasificación de módulos
finitamente generados sobre un anillo de valoración discreta.
Estas $r$ filas generan $R$ y son linealmente independientes.

Tómense los elementos $g_i$ de cuyas matrices proceden esas
filas, suprimiendo repeticiones. Las filas de $B$ contienen la
base elegida y todas pertenecen a $R$; generan exactamente $R$.
Ser anulado por ellas equivale a ser anulado por cada
$\rho(g)-I_d$, que es la condición de invariancia.
\end{proof}

Los $g_i$ no tienen por qué generar $G$. El lema reduce un sistema
de ecuaciones lineales, no determina todos los cociclos del grupo
por sus valores en esos elementos. Tampoco asegura una cota para
las valoraciones de las filas escogidas.

\begin{proposicion}[Smith y certificado de longitud]
\label{prop:lu-smith}
Sean $\varpi^{a_1},\ldots,\varpi^{a_r}$ los divisores elementales
no nulos de $B$, con $0\le a_1\le\cdots\le a_r$. Existe una
isomorfía, no necesariamente canónica,
\begin{equation}
 A_d^G\simeq\bigoplus_{i=1}^r\mathcal O/\varpi^{a_i}\mathcal O
                    \ \oplus\ (K/\mathcal O)^{d-r}.
 \label{eq:lu-smith}
\end{equation}
Son equivalentes la longitud finita de $A_d^G$, la igualdad $r=d$
y la anulación de $(K^d)^G$. En ese caso,
\begin{equation}
 \length_{\mathcal O}A_d^G
   =\sum_{i=1}^d a_i
   =v_\varpi(I_d(B))
   =\min_Jv_\varpi(\det B_J),
 \label{eq:lu-menores}
\end{equation}
donde $I_d(B)$ es el ideal de los menores de tamaño $d$.
Para $d=0$ se toma el menor vacío igual a uno.
\end{proposicion}
\begin{proof}
Las operaciones invertibles que llevan $B$ a su forma de Smith
inducen automorfismos en los módulos divisibles. El núcleo de
multiplicar por $\varpi^{a_i}$ en $K/\mathcal O$ es
$\varpi^{-a_i}\mathcal O/\mathcal O\simeq\mathcal O/\varpi^{a_i}$.
Las $d-r$ coordenadas restantes no tienen restricción; esto prueba
\eqref{eq:lu-smith}. Sobre $K$, el mismo sistema tiene núcleo de
dimensión $d-r$, lo que prueba las equivalencias.

Cuando $r=d$, la longitud es la suma de los exponentes. Las
operaciones invertibles conservan el ideal de menores: la fórmula
de Cauchy--Binet da una inclusión y las operaciones inversas dan
la otra. En la diagonal, ese ideal es
$(\varpi^{\sum_i a_i})$. La valoración de un ideal finitamente
generado sobre $\mathcal O$ es el mínimo de las valoraciones de
sus generadores. De aquí resulta \eqref{eq:lu-menores}.
\end{proof}

La fórmula proporciona un certificado numérico: una cota
$\lambda_{\mathcal O}(A_d^G)\le C$ equivale a la existencia de
un menor de tamaño $d$ formado por filas de los operadores
$\rho(g)-I_d$ cuya valoración sea a lo sumo $e_{\mathcal O}C$.
En efecto, las filas originales generan $R$, por lo que sus
menores generan el mismo ideal de determinantes; en un anillo
local alguno realiza la valoración mínima. Si $r<d$, el módulo
contiene una parte divisible no nula y no existe tal certificado
de longitud finita. El ideal de menores y su independencia de
presentación son también el mecanismo de Fitting de la
Nota~\ref{sec:theta-integral}; véase Stacks, etiqueta 07Z6.

Puede ser imposible detectar los invariantes con una sola matriz.
Para $G=C_2\times C_2$ sobre $\mathcal O=\Z_2$, tome
\begin{equation}
 \rho(a)=\operatorname{diag}(1,-1,-1),\qquad
 \rho(b)=\operatorname{diag}(-1,1,-1).
 \label{eq:lu-klein}
\end{equation}
Cada elemento de $G$ tiene un vector fijo racional no nulo;
por tanto $\det(\rho(g)-I_3)=0$ para todo $g$. Sin embargo, las
filas conjuntas generan $2\Z_2^3$, y los invariantes de
$(\Q_2/\Z_2)^3$ son $(\Z/2\Z)^3$, de longitud tres.

\subsection{Cuánto se pierde al pasar de coeficientes finitos a divisibles}
\label{subsec:lu-bockstein}

\begin{proposicion}[Longitud exacta del núcleo de coeficientes]
\label{prop:lu-bockstein}
Con la notación de \eqref{eq:lu-smith}, para $n\ge0$ se tiene
\begin{equation}
 \ker\!\left[H^1(G,A_d[\varpi^n])\to H^1(G,A_d)\right]
       \simeq A_d^G/\varpi^n A_d^G,
 \label{eq:lu-bockstein}
\end{equation}
y este núcleo tiene longitud finita igual a
\begin{equation}
 \sum_{i=1}^r\min(n,a_i)\le rn\le dn.
 \label{eq:lu-bockstein-longitud}
\end{equation}
La afirmación sigue siendo válida si $A_d^G$ tiene longitud infinita.
\end{proposicion}
\begin{proof}
Dado $b\in A_d^G$, elegimos $a\in A_d$ con $\varpi^n a=b$.
El cociclo $g\mapsto ga-a$ es continuo y toma valores en
$A_d[\varpi^n]$. Otra preimagen difiere de $a$ en un elemento
de ese submódulo, de modo que define la misma clase.
La clase es cero precisamente cuando existe
$c\in A_d[\varpi^n]$ con $ga-a=gc-c$ para todo $g$; esto
equivale a $a-c\in A_d^G$, y luego a
$b\in\varpi^n A_d^G$. Recíprocamente, todo cociclo del núcleo
es un borde $ga-a$ con valores en $A_d[\varpi^n]$, de donde
$\varpi^n a\in A_d^G$. Queda probado \eqref{eq:lu-bockstein}.

En \eqref{eq:lu-smith}, cada sumando divisible coincide con su
múltiplo por $\varpi^n$, y no aporta al cociente. Cada sumando
finito aporta longitud $\min(n,a_i)$. La suma da la fórmula
y, al tener $r$ sumandos, las dos desigualdades.
\end{proof}

Para un nivel $\ell^q$ se toma $n=e_{\mathcal O}q$, y así
\begin{equation}
 \lambda_{\mathcal O}\ker\!\left[
 H^1(G,A_d[\ell^q])\to H^1(G,A_d)\right]\le dq.
 \label{eq:lu-bockstein-normalizado}
\end{equation}
La sucesión de la Nota~\ref{sec:construccion-selmer} queda ahora
acompañada de una medida exacta, incluso con invariantes infinitos.
Para una representación fija, al crecer $n$ la longitud se
estabiliza en $\sum_i a_i$; no se estabiliza necesariamente en
la longitud de todos los invariantes.

Las extensiones integrales importan. Tome $K=\Q_\ell$ y
$\mathcal O=\Z_\ell$. Para $G=\Z_\ell$ y $N\ge1$,
considere la representación continua
\begin{equation}
 \rho_N(x)=\begin{pmatrix}1&\ell^N x\\0&1\end{pmatrix}.
 \label{eq:lu-unipotente}
\end{equation}
Sus invariantes divisibles son
$(\Q_\ell/\Z_\ell)\oplus(\Z_\ell/\ell^N)$; por tanto la
longitud del núcleo al nivel $\ell^q$ es $\min(q,N)$.
La semisimplificación es trivial de rango dos, cuyos invariantes
son completamente divisibles y cuyo núcleo es cero.
Conocer solo los caracteres de la semisimplificación pierde
justamente la parte finita que mide \eqref{eq:lu-bockstein}.

\subsection{Una cota explícita para la ventana de torcidos negativos}
\label{subsec:lu-h0-caracteres}

Volvamos al cuerpo de números $E$, con $n_E=[E:\Q]$, y supongamos
que $\mathcal C$ es un conjunto finito de lugares sobre $\ell$.
Para $w\mid\ell$, sea $\epsilon_w=e(E_w/\Q_\ell)$. El
Lema~\ref{lem:fu-invariantes} ya demuestra, para $e\ge1$,
\begin{equation}
 a_{\ell,e}(E_w):=\length_{\Z_\ell}
 H^0(E_w,\Q_\ell/\Z_\ell(-e))
 \le s_\ell+v_\ell(\epsilon_w)+v_\ell(e),\qquad
 s_\ell=\begin{cases}1,&\ell\ne2,\\2,&\ell=2.\end{cases}
 \label{eq:lu-ciclotomica-previa}
\end{equation}
Aplicamos ese resultado al supremo concreto de la construcción.

\begin{proposicion}[Control del supremo de invariantes negativos]
\label{prop:lu-h0}
Para $t\ge2$, el número $H_0$ de \eqref{eq:lu-h0-def} es finito y
\begin{equation}
 \frac{H_0}{e_{\mathcal O}}
 \le s_\ell+\max_{w\in\mathcal C}v_\ell(\epsilon_w)
                 +\lfloor\log_\ell(t-1)\rfloor
 \le s_\ell+\lfloor\log_\ell n_E\rfloor
                 +\lfloor\log_\ell(t-1)\rfloor.
 \label{eq:lu-h0-cota}
\end{equation}
La primera expresión se usa para $\mathcal C\ne\varnothing$;
si $\mathcal C=\varnothing$ o $t\le1$, entonces $H_0=0$.
Además, $H_0=0$ si $\ell>(t-1)n_E+1$.
\end{proposicion}
\begin{proof}
La comparación de coeficientes de la
Proposición~\ref{prop:ti-escalares}, en grado cero, da
\[
 \length_{\mathcal O}H^0(E_w,K/\mathcal O(-e))
       =e_{\mathcal O}a_{\ell,e}(E_w).
\]
En el intervalo $1\le e<t$, el máximo de $v_\ell(e)$ es
$\lfloor\log_\ell(t-1)\rfloor$: ninguna potencia mayor cabe en
el intervalo y la potencia correspondiente sí cabe. Por otra
parte $\epsilon_w\le[E_w:\Q_\ell]\le n_E$, de modo que
$v_\ell(\epsilon_w)\le\lfloor\log_\ell n_E\rfloor$.
Tomar el máximo finito en \eqref{eq:lu-ciclotomica-previa} prueba
la finitud y la cota.

Si $\ell>(t-1)n_E+1$ y el conjunto de índices es no vacío,
entonces $\ell$ es impar. La condición de soporte del
Lema~\ref{lem:fu-invariantes} exige
$\ell-1\mid e\epsilon_w$ para un invariante no nulo.
Pero $0<e\epsilon_w\le(t-1)n_E<\ell-1$, lo que es imposible.
\end{proof}

Con $n_E\le N_*$ y $2\le t\le T_*$, la cantidad
\[
 B_{\chi,*}:=2+\lfloor\log_2N_*\rfloor
                    +\lfloor\log_2(T_*-1)\rfloor
\]
acota $H_0/e_{\mathcal O}$ para todos los primos, y los primos
que pueden contribuir satisfacen
$\ell\le(T_*-1)N_*+1$. Son conclusiones uniformes para este
proveedor de caracteres. Si el grado o el corte crecen, esas
mismas fórmulas muestran su dependencia. No controlan el bloque
trivial, el bloque automorfo ni los demás términos de Selmer.

Cuando la imagen ciclotómica local es exactamente
$1+\ell^b\Z_\ell$, con $b\ge1$ si $\ell$ es impar y $b\ge2$
si $\ell=2$, el cálculo de la
Proposición~\ref{prop:ts-invariantes} da
\begin{equation}
 \sum_{e=1}^{t-1}a_{\ell,e}(E_w)
       =(t-1)b+v_\ell((t-1)!).
 \label{eq:lu-suma-caracteres}
\end{equation}
Es una información más fina que el máximo. Sin embargo, no
permite sustituir sin prueba el término $t|\mathcal C|H_0$ de
\eqref{eq:ti-errores-explicitos} por esa suma: tal sustitución
requiere rehacer la comparación que produjo el factor $t$.

\subsection{Subcocientes divisibles con exponente acotado}
\label{subsec:lu-subcocientes}

\begin{lema}[Exponente y corango controlan la longitud]
\label{lem:lu-subcocientes}
Sean $C\subseteq B\subseteq(K/\mathcal O)^D$ submódulos y
$Q=B/C$. Si $\varpi^hQ=0$, con $h\ge0$, entonces $Q$ tiene
longitud finita y
\begin{equation}
 \length_{\mathcal O}Q\le Dh.
 \label{eq:lu-subcocientes}
\end{equation}
No se supone que $B$, $C$ o $Q$ sean finitamente generados.
\end{lema}
\begin{proof}
Tomemos cualquier submódulo finitamente generado $Q_0$ de $Q$.
Levantamos sus generadores a elementos $b_1,\ldots,b_m$ de $B$
y ponemos $R=\sum_i\mathcal O b_i$. Cada $b_i$ es torsión;
por tanto un mismo $\varpi^N$ anula $R$, que queda contenido en
$(K/\mathcal O)^D[\varpi^N]\simeq(\mathcal O/\varpi^N)^D$.
Así, $R$ tiene longitud finita.

Por Smith, el número mínimo de generadores de $R$ es
$\dim_k R[\varpi]$. Como $R[\varpi]$ se incluye en $k^D$, ese
número es a lo sumo $D$. Su imagen $Q_0$ tiene, en consecuencia,
a lo sumo $D$ generadores, y la hipótesis de exponente da una
sobreyección $(\mathcal O/\varpi^h)^D\twoheadrightarrow Q_0$.
De aquí $\length Q_0\le Dh$.

Todas las longitudes de submódulos finitamente generados de $Q$
son enteros entre cero y $Dh$. Elija uno de longitud máxima.
Si un elemento de $Q$ quedara fuera, añadirlo produciría un
submódulo finitamente generado estrictamente mayor, y la
aditividad de longitudes daría una longitud mayor. Esto es
imposible. El submódulo elegido es $Q$ entero, lo que prueba
su finitud y \eqref{eq:lu-subcocientes}.
\end{proof}

La prueba evita inferir finitud a partir del exponente solamente:
$(\mathcal O/\varpi)^{(\N)}$ tiene exponente $\varpi$ y longitud
infinita. El control del módulo ambiente es la otra hipótesis
esencial. Si el ambiente ya es
$M=\bigoplus_i\mathcal O/\varpi^{a_i}$, la cota se afina a
$\sum_i\min(h,a_i)$: $Q$ es cociente de $B/\varpi^hB$,
cuya longitud es $\length B[\varpi^h]$, y
$B[\varpi^h]\subseteq M[\varpi^h]$.

\subsection{Cambio de base aritmética con pérdida explícita}
\label{subsec:lu-restriccion}

\begin{teorema}[Núcleo de restricción a un subgrupo normal abierto]
\label{thm:lu-restriccion}
Sea $H\triangleleft G$ abierto, sea $\Gamma=G/H$ de orden $n$,
y supongamos que $\Gamma$ puede generarse por $s$ elementos.
Para $A_d$ como antes, sin exigir finitud de $A_d^H$, se tiene
\begin{equation}
 \length_{\mathcal O}\ker\!\left[
 H^1(G,A_d)\xrightarrow{\operatorname{res}}H^1(H,A_d)\right]
       \le d s e_{\mathcal O}v_\ell(n).
 \label{eq:lu-restriccion}
\end{equation}
En particular, el núcleo es finito y la restricción es inyectiva
si $\ell\nmid n$.
\end{teorema}
\begin{proof}
Pongamos $M=A_d^H$. Primero identificamos el núcleo con
$H^1(\Gamma,M)$. Una clase cuya restricción es cero admite,
tras restar un coborde, un cociclo $z$ que se anula en $H$.
La identidad de cociclo muestra $z(gh)=z(g)$ para $h\in H$.
Además $hg=g(g^{-1}hg)$ y la normalidad dan
$hz(g)=z(hg)=z(g)$, por lo que los valores pertenecen a $M$.
Así $z$ desciende a $\Gamma$. A la inversa se puede inflar todo
cociclo de $\Gamma$. Si su inflación es un borde $ga-a$, el
hecho de que se anule en $H$ implica $a\in M$; era ya un borde
de $\Gamma$ con valores en $M$. Los cociclos descendidos son
continuos porque $\Gamma$ es finito y $H$ abierto.

Un cociclo de $\Gamma$ queda determinado por sus valores en los
$s$ generadores: se usan sucesivamente
$z(uv)=z(u)+uz(v)$ y $z(u^{-1})=-u^{-1}z(u)$.
Se obtiene una inyección
\[
 Z^1(\Gamma,M)\hookrightarrow M^s\subseteq(K/\mathcal O)^{ds}.
\]
Por tanto $H^1(\Gamma,M)$ es un subcociente de ese módulo.
Para cada cociclo $z$, sea $b=\sum_{\gamma\in\Gamma}z(\gamma)$.
La permutación $\gamma\mapsto g\gamma$ da
$b=nz(g)+gb$, es decir $nz(g)=b-gb$. El entero $n$ anula todas
las clases. En $\mathcal O$ es una unidad por
$\varpi^{e_{\mathcal O}v_\ell(n)}$. El
Lema~\ref{lem:lu-subcocientes}, con $D=ds$, termina la prueba.
\end{proof}

Puede tomarse $s\le\lfloor\log_2 n\rfloor$: al añadir un
generador exterior, el orden del subgrupo generado se multiplica
al menos por dos. Si $n=1$, se toma $s=0$ y el núcleo es cero.
Para una extensión galoisiana de cuerpos aritméticos $L/E$, la
normalización conserva la dependencia en su grado:
\begin{equation}
 \lambda_{\mathcal O}(\ker\operatorname{res}_{L/E})
       \le d s v_\ell([L:E]).
 \label{eq:lu-restriccion-normalizada}
\end{equation}
El factor $e_{\mathcal O}$ desaparece; el grado aritmético
permanece. La inyectividad especial de las capas ciclotómicas
negativas demostrada en la Nota~\ref{sec:transporte-selmer}
es más fuerte en esa situación y se conserva.

La dependencia puede ser efectiva. Para $G=\Z_\ell$,
$H=\ell^a\Z_\ell$ y coeficientes triviales
$A=\Q_\ell/\Z_\ell$, la restricción es la multiplicación por
$\ell^a$ al identificar un homomorfismo continuo con su valor
en un generador. Todo elemento de $A$ tiene orden finito, así
que efectivamente determina un homomorfismo continuo. El núcleo
tiene longitud $a$ y alcanza \eqref{eq:lu-restriccion-normalizada}
con $d=s=1$.

\subsection{Qué condiciones permiten transportar Selmer}
\label{subsec:lu-selmer}

Sea $L/E$ una extensión finita galoisiana y sea $S$ un conjunto
finito de lugares que contiene los de ramificación de la
extensión y de la representación, los lugares sobre $\ell$ y
los infinitos. Se utiliza el mismo problema de ramificación
para $G_{E,S}$ y $G_{L,S_L}$. Si el problema inicial solo permitía
ramificación sobre $\ell$, ampliar $S$ cambia el grupo global
y necesita una comparación adicional.

Dadas condiciones locales $F_v\subseteq H^1(E_v,A_d)$ y
$F_w\subseteq H^1(L_w,A_d)$, definimos
\[
 S_F(E)=\ker\!\left[H^1(G_{E,S},A_d)\to
             \prod_{v\in S}H^1(E_v,A_d)/F_v\right],
\]
y análogamente $S_F(L)$. Los mapas de localización se fijan con
encajes compatibles y las identificaciones por conjugación
necesarias.

\begin{proposicion}[Transporte de longitud con condiciones locales]
\label{prop:lu-selmer}
Supongamos, para todo $w\mid v$, que
\begin{equation}
 \operatorname{res}_{L_w/E_v}(F_v)\subseteq F_w.
 \label{eq:lu-compatibilidad-local}
\end{equation}
Si $S_F(L)$ tiene longitud finita y $\operatorname{Gal}(L/E)$
admite $s$ generadores, entonces $S_F(E)$ tiene longitud finita y
\begin{equation}
 \lambda_{\mathcal O}S_F(E)
 \le\lambda_{\mathcal O}S_F(L)+d s v_\ell([L:E]).
 \label{eq:lu-selmer}
\end{equation}
\end{proposicion}
\begin{proof}
Las restricciones global y local forman un cuadrado conmutativo
para cada $w\mid v$. La hipótesis
\eqref{eq:lu-compatibilidad-local} hace que la restricción global
lleve $S_F(E)$ a $S_F(L)$. Su núcleo se incluye en el núcleo de
restricción de toda la cohomología, cuya longitud acota el
Teorema~\ref{thm:lu-restriccion}. La imagen es un submódulo de
$S_F(L)$. La sucesión exacta de núcleo e imagen demuestra la
finitud del dominio y, por aditividad, la desigualdad.
\end{proof}

Para condiciones locales cero, la inclusión requerida es
automática. Para condiciones de periodos debe demostrarse su
functorialidad bajo la extensión utilizada; no basta que ambos
submódulos reciban el nombre de ``condición finita''.
Si se comparan en un mismo cuerpo condiciones
$F_v\subseteq F'_v$, la localización induce una inyección
\begin{equation}
 S_{F'}(E)/S_F(E)\hookrightarrow\bigoplus_{v\in S}F'_v/F_v.
 \label{eq:lu-cambio-condicion}
\end{equation}
El núcleo es exactamente $S_F(E)$, por definición. Así se puede
sumar el coste de relajar las condiciones cuando esos cocientes
son finitos; no se afirma sobreyectividad sobre la suma local.

Tampoco se deduce sobreyectividad de una restricción de índice
primo con $\ell$. Para $\ell$ impar, tome
$G=\Z_\ell\rtimes C_2$, con acción por inversión, y coeficientes
triviales $A=\Q_\ell/\Z_\ell$. Todo homomorfismo $G\to A$
se anula en $C_2$ y debe ser igual a su negativo en $\Z_\ell$;
como $2$ es una unidad, es cero. Pero
$H^1(\Z_\ell,A)=A$. La restricción tiene núcleo cero y
conúcleo infinito. La cota anterior solo usa su núcleo.

\subsection{Insertar los controles en el presupuesto de la Nota 25}
\label{subsec:lu-presupuesto}

Escribamos
\[
 D_\delta=\frac{v_\delta}{e_{\mathcal O}},\qquad
 B_0=\max_{w\in\mathcal C}\frac{h_{0,w}}{e_{\mathcal O}},\qquad
 B_\chi=\frac{H_0}{e_{\mathcal O}}.
\]
El máximo vacío es cero. La Nota~\ref{sec:theta-integral}
conserva por separado las filas, el bloque automorfo, los Selmer
negativos y el grupo de clases. Ahora podemos nombrar un
certificado para cada error local.

\begin{proposicion}[Presupuesto uniforme a partir de datos acotados]
\label{prop:lu-presupuesto}
Considérese una familia para la que se cumplen las hipótesis de
\eqref{eq:ti-local-display}. Supongamos que, antes de elegir
sus miembros, se han fijado cotas no negativas
\begin{equation}
\begin{gathered}
 |\mathcal C|\le C_*,\quad j_w\le J_*,\quad e'_w\le E_*,\quad
 D_\delta\le D_*,\quad B_0\le B_{0,*},\quad B_\chi\le B_{\chi,*},\\
 t\le T_*,\quad n_0\le N_{0,*},\quad c_B\le C_{B,*},\quad
 \sum_{w\in\mathcal R}(d_w+2)\le R_*,\quad
 \sum_{w\in\mathcal Q}\frac{v_{\zeta,w}}{e_{\mathcal O}}\le Z_*.
\end{gathered}
\label{eq:lu-datos-uniformes}
\end{equation}
Entonces los errores de \eqref{eq:ti-errores-explicitos}
satisfacen simultáneamente
\begin{equation}
\begin{aligned}
 \frac{\mathcal E_{\rm col}}{e_{\mathcal O}}
  &\le C_*\bigl[(1+J_*)E_*+(J_*+2)D_*+B_{0,*}\bigr],\\
 \frac{\mathcal E_{\rm loc}}{e_{\mathcal O}}
  &\le T_*C_{B,*}R_*+4(N_{0,*}+T_*)(C_*+1)D_*\\
  &\hspace{1.5em}+T_*C_*B_{\chi,*}+T_*Z_*.
\end{aligned}
\label{eq:lu-presupuesto-uniforme}
\end{equation}
\end{proposicion}
\begin{proof}
Después de dividir la primera igualdad de
\eqref{eq:ti-errores-explicitos} por $e_{\mathcal O}$, cada
sumando es
$(1+j_w)e'_w+(j_w+2)D_\delta+h_{0,w}/e_{\mathcal O}$.
Todos los factores son no negativos. Sustituirlos por sus cotas
y usar $|\mathcal C|\le C_*$ da la primera desigualdad.
La misma sustitución, término a término en la segunda igualdad,
da la segunda. No se ha cambiado la multiplicidad de ningún
bloque ni se ha absorbido un coste sin exhibir su dependencia.
\end{proof}

Hay avances efectivos dentro de estas hipótesis. Para columnas
con las ventanas y entradas de periodos indicadas, la
Proposición~\ref{prop:lu-ventana-comun} construye $E_*$.
Para los caracteres negativos, la
Proposición~\ref{prop:lu-h0} construye $B_{\chi,*}$ a partir
del grado y del corte. Para el bloque automorfo, una
identificación integral equivariante con $A_d$ y menores de
valoración normalizada acotada proporcionan $B_{0,*}$ mediante
\eqref{eq:lu-h0-bloque} y \eqref{eq:lu-menores}.
El Teorema~\ref{thm:lu-restriccion} acota el coste de transportar
los grupos globales a una base galoisiana de grado controlado.

Las filas conservan su cota propia. Si se construye un
levantamiento en esa clase, con rangos $r_w$, exponentes $a_w$
y torsiones $T_w$, \eqref{eq:pl-coste-fino} da
\[
 \frac{L_{{\rm row},w}}{e_{\mathcal O}}
       \le r_w a_w+\frac{\length T_w}{e_{\mathcal O}}.
\]
La otra ruta es la comparación de periodos
\eqref{eq:pl-colle}--\eqref{eq:pl-tie}. Ninguna permite reutilizar
sin prueba un exponente calculado solo para las columnas.
Los términos globales $s_\rho$, los Selmer negativos y el grupo
de clases de \eqref{eq:ti-local-display} mantienen sus propias
cotas. Una cota uniforme por primo tampoco demuestra por sí
sola una suma ponderada sobre todos los primos; ese paso sigue
exigiendo las hipótesis del Lema~\ref{lem:ti-presupuesto}.

\subsection{Ejercicios resueltos}

\begin{ejercicio}[Corregir el representante sin pagar otro exponente]
\label{ex:lu-1}
Sea $D$ un grupo, $p:M_1\twoheadrightarrow M_2$ equivarante,
y sean $K_0,\kappa$ cociclos con
$pK_0-\ell^e\kappa=dv$. Elija $w$ con $p(w)=v$.
Compruebe la corrección, su signo y la composición de dos
levantamientos de exponentes $e_1,e_2$.
\end{ejercicio}
\paragraph{Solución.}
Se toma $K=K_0-dw$, pues $p(dw)=dv$; con el signo positivo
quedaría $2dv$ en la diferencia. La clase de $K$ es la de $K_0$.
Si $M_1\to M_2\to M_3$ son sobreyecciones compatibles con las
condiciones locales, se levanta primero
$\ell^{e_2}\kappa$ a $K_2$ en $M_2$ y luego
$\ell^{e_1}K_2$ a $K_1$ en $M_1$. La composición envía
$K_1$ a $\ell^{e_1+e_2}\kappa$. Los exponentes se suman;
no se multiplican. Para una lista finita de cociclos se repite
la construcción con el mismo exponente: el número de cociclos
puede costar longitud después, pero no añade exponente aquí.

\begin{ejercicio}[El coste de torsión de una fila]
\label{ex:lu-2}
Sea $e_{\mathcal O}=2$,
$F=\mathcal O^2\oplus\mathcal O/\varpi^5$ e $I=\ell^2F$.
Calcule $\length(F/I)$ y compárela con la cota por rango y
torsión de \eqref{eq:pl-coste-fino}.
\end{ejercicio}
\paragraph{Solución.}
$\ell^2$ es una unidad por $\varpi^4$. Los dos sumandos libres
aportan $4+4$, y el sumando de torsión aporta $\min(5,4)=4$.
La longitud exacta es $12$. La cota que conserva toda la torsión
da $2\cdot2\cdot2+5=13$. Usar solo el rango daría $8$ y
fallaría. Las longitudes normalizadas exacta y superior son
respectivamente $6$ y $13/2$.

\begin{ejercicio}[Menores y pérdida al nivel finito]
\label{ex:lu-3}
Suponga que las filas de un sistema de invariancia generan el
mismo submódulo que
\[
 B=\begin{pmatrix}\varpi^2&0\\0&\varpi^5\\
                       \varpi^3&\varpi^6\end{pmatrix}.
\]
Calcule la longitud de los invariantes y la del núcleo de
\eqref{eq:lu-bockstein} para $n=3$. Explique qué sucede en el
ejemplo de Klein \eqref{eq:lu-klein} para un primo impar.
\end{ejercicio}
\paragraph{Solución.}
La tercera fila es $\varpi$ veces la suma de las dos primeras.
El ideal de entradas es $(\varpi^2)$ y el de menores de tamaño
dos es $(\varpi^7)$; por tanto los exponentes de Smith son
$2$ y $5$. Los invariantes tienen longitud $7$, y el núcleo
finito tiene longitud $\min(3,2)+\min(3,5)=5$.
En el ejemplo de Klein, las filas conjuntas generan
$2\mathcal O^3$. Si $\ell$ es impar, $2$ es una unidad y los
invariantes divisibles son cero, aunque cada matriz individual
siga siendo singular sobre $K$.

\begin{ejercicio}[La semisimplificación no conserva el núcleo]
\label{ex:lu-4}
Para \eqref{eq:lu-unipotente}, calcule los invariantes y la
pérdida al nivel $\ell^q$ cuando $N=4$, $q=6$. Compare con la
representación trivial de rango dos.
\end{ejercicio}
\paragraph{Solución.}
Un par $(u,v)$ es fijo si y solo si $\ell^4v=0$; basta probarlo
en $x=1$ y entonces vale para todo $x\in\Z_\ell$.
Así $u$ es arbitrario en $\Q_\ell/\Z_\ell$ y $v$ pertenece
a su subgrupo de orden $\ell^4$. La parte divisible se anula
en el cociente por $\ell^6$, mientras la parte finita aporta
longitud $4$. En la representación trivial ambos sumandos son
divisibles y la pérdida es cero. Las dos representaciones tienen
la misma semisimplificación racional.

\begin{ejercicio}[Un núcleo de restricción que alcanza la cota]
\label{ex:lu-5}
Tome coeficientes triviales $A=\Q_\ell/\Z_\ell$,
$G=\Z_\ell$ y $H=\ell^3\Z_\ell$. Calcule el núcleo de
restricción y explique por qué no se necesita que $A^H$ sea finito.
\end{ejercicio}
\paragraph{Solución.}
Un cociclo es un homomorfismo continuo, determinado por su valor
en $1$. Al restringir, el generador elegido de $H$ es $\ell^3$,
por lo que el mapa es $a\mapsto\ell^3a$ en $A$. El núcleo es
$A[\ell^3]\simeq\Z/\ell^3$, de longitud $3$.
El cociente $G/H$ es cíclico de orden $\ell^3$, y
$d=s=e_{\mathcal O}=1$; la cota es exactamente $3$.
Aunque $A^H=A$ sea infinito, sus cociclos sobre el cociente finito
producen un subcociente de corango uno anulado por $\ell^3$.
Es precisamente el caso cubierto por el lema de subcocientes.

\begin{ejercicio}[Un máximo ciclotómico y una ventana de columnas]
\label{ex:lu-6}
Sea $L=\Q_3(\mu_3)$, cuya imagen ciclotómica es $1+3\Z_3$,
y tome $t=6$. Calcule las cinco longitudes de invariantes
negativos, su máximo y su suma. Por separado, sobre una base local
no ramificada, suponga que las hipótesis del productor de columnas
se cumplen con $C_A=2$ y $M+h\le12$. ¿Qué exponente común
proporciona \eqref{eq:lu-exponente-comun}?
\end{ejercicio}
\paragraph{Solución.}
Para $e=1,2,3,4,5$, las longitudes son
$1+v_3(e)$, es decir $1,1,2,1,1$. Su máximo es $2$ y su suma
es $6=5+v_3(5!)$. Como $[L:\Q_3]=2$, la cota local por grado
también da $1+\lfloor\log_3 2\rfloor+\lfloor\log_3 5\rfloor=2$.
En la segunda aplicación, el exponente de columnas puede elegirse como
\[
 E_*=\left\lceil2\bigl(2+12/3+\log_3 13\bigr)\right\rceil=17.
\]
El cuerpo $L$ del primer cálculo es ramificado y no se utiliza
como base del productor de la segunda aplicación.
El máximo $2$ mide
invariantes de caracteres; $17$ es un exponente permitido por
el productor de periodos bajo sus hipótesis. Ninguno determina
por sí solo la longitud de una fila ni una constante para toda
la torre.

\subsection{Alcance alcanzado y siguiente construcción}
\label{subsec:lu-frontera}

La nota desarrolla trece resultados con demostración y seis
ejercicios resueltos. La corrección de cociclos y la selección
finita permiten construir las funciones de columnas desde la
entrada de levantamiento. La presentación de filas conserva los
bordes integrales, la topología y la torsión. Los invariantes
reciben un certificado por menores; el paso de coeficientes y el
cambio de base aritmética reciben cotas de núcleo demostradas.
El supremo de caracteres negativos queda acotado por grado y
corte, con soporte primo explícito cuando ambos están acotados.

La concordancia exacta, los hashes y los rangos leídos se conservan
en \texttt{fuentes/README.md} y \texttt{REGISTRO.json}. La lectura
de las entradas \texttt{LiftTheoremAllCases},
\texttt{LengthBoundAtPlaces} y \texttt{H0.FinCore} se distingue
de una prueba de todos sus productores. No se ejecutó Lean ni
se inspeccionó su clausura transitiva. A93, pp.~62--64, conserva
la misma exigencia de construir los mapas, condiciones locales
y multiplicidades antes de sumar longitudes; esta etapa no
constituye una certificación global de la propuesta.

El siguiente desarrollo concreto es realizar esos certificados
en la familia theta variable: una ventana de saltos compatible
con sus bases locales, el retículo matricial y las condiciones
de periodos de las filas, y menores de valoración controlada en
el bloque automorfo. Deben acompañarse de cotas para los
separadores, los generadores, los grados y las pérdidas auxiliares.
El presupuesto \eqref{eq:lu-presupuesto-uniforme} indica dónde
entra cada una y qué dependencia permanece cuando crece la torre.
